\chapter{Guide de Déploiement et d'Exécution Locale}

\section{Prérequis Logiciels}

Pour compiler et exécuter le prototype en environnement local, les outils suivants sont requis :
\begin{itemize}
    \item \textbf{SDK .NET 9.0} (ou version supérieure) ;
    \item \textbf{Node.js} (version 20.x LTS ou 22.x) et \textbf{npm} (version 10.x+) ;
    \item \textbf{Angular CLI 19} (\texttt{npm install -g @angular/cli}) ;
    \item \textbf{Microsoft SQL Server} ou \textbf{LocalDB} (inclus avec Visual Studio).
\end{itemize}

\section{Procédure de Lancement Pas-à-Pas}

\subsection{1. Démarrage de l'Application Frontend Angular}
Ouvrir un premier terminal et exécuter :
\begin{lstlisting}[language=bash]
cd "Digital twin (projet OCP)/frontend"
npm install
npm start
\end{lstlisting}
L'interface de supervision est accessible à l'adresse : \url{http://localhost:4200}.

\subsection{2. Lancement du Serveur API Backend}
Ouvrir un deuxième terminal et exécuter :
\begin{lstlisting}[language=bash]
cd "Digital twin (projet OCP)/backend/DigitalTwin.Api"
dotnet run
\end{lstlisting}
L'API REST et le Hub SignalR démarrent sur : \url{https://localhost:7123} (Swagger : \url{https://localhost:7123/swagger}).

\subsection{3. Démarrage du Simulateur d'Usine}
Ouvrir un troisième terminal et exécuter :
\begin{lstlisting}[language=bash]
cd "Digital twin (projet OCP)/backend/DigitalTwin.Simulator"
dotnet run
\end{lstlisting}
Le simulateur injecte immédiatement les flux de télémétrie et rejoue le cycle de vie nominal et d'incidents sur les équipements.

\subsection{4. Exécution des Tests Automatisés}
Pour exécuter la suite des 83 tests unitaires du backend :
\begin{lstlisting}[language=bash]
cd "Digital twin (projet OCP)/backend"
dotnet test
\end{lstlisting}
